\documentclass[../main.tex]{subfiles}
\begin{document}
%==========================================TIÊU ĐỀ TẬP========================================
\begin{table}[h]
  \begin{adjustwidth}{-1cm}{}
    \begin{tabular}{>{\centering\arraybackslash}p{6cm}|>{\raggedright\arraybackslash}p{12.5cm}}
      \multirow{5}{6cm}
      % Cột bên trái: ÉPISODE và số tập
      {\\[-40pt]
        \flaregothic\fontsize{20pt}{20pt}\selectfont \centering
        {\contour{errorfive}{\textcolor{black}{ÉPISODE}}}\color{black}\\
        \burbank\fontsize{72pt}{72pt}\selectfont
        \includegraphics[height=3.12359cm]{../../../../Image/Book?/number/num25.png}
        \\[18pt]
        % \flaregothic\fontsize{14pt}{14pt}\selectfont
        % \color{epattr}BIENVENUE, \uline{\color{Banpire} Mel} !
        % \flaregothic\fontsize{18pt}{18pt}\selectfont
        % \color{epattr}ÉPISODE PERDU
        \color{black}
      }
       &
      % Dòng thứ nhất: tên tập tiếng Nhật
      {\fotskip\fontsize{18pt}{18pt}\selectfont \ruby{子}{こ}\ruby{供}{ども}の\ruby{霊}{れい}}                         \\[6pt]
      % Dòng thứ hai: tên tập tiếng Anh
       & {\cafeteria\fontsize{18pt}{18pt}\selectfont Child Spirits}                                           \\[4pt]
      % Dòng thứ ba: tên tập tiếng Việt
       & {\vnmsans\fontsize{18pt}{18pt}\selectfont Hồn ma trẻ em}                                             \\[8pt]
      % Dòng thứ tư: Nhân vật xuất hiện
       & {\caviar{\fontsize{14pt}{14pt}\selectfont Nhân vật: \emp{sora} \emp{miko} \emp{korone} \emp{okayu}}}
      % \\[6pt]
      % % Dòng cuối cùng: Thời gian diễn ra của tập (thường sẽ là ngày phát hành)
      % & {\continuum\fontsize{16pt}{16pt}\selectfont Ngày phát hành: 11/06/2022}\\[8pt]
    \end{tabular}
  \end{adjustwidth}
\end{table}
%===========================================NỘI DUNG==========================================
\color{black}

\pov{Hikawa Kuon}

Ngày đầu tiên sau khi thế giới được tái sinh\dots\ nghe thì có vẻ hoành tráng lắm, nhưng buổi học sáng của tôi lại bắt đầu y như những ngày bình thường của một học sinh cấp ba tiêu chuẩn: tập trung, ghi chép, nghe giảng, thỉnh thoảng quay sang nhìn Kaigou-chan đang hí hoáy với cuốn sổ màu đỏ, còn Suzuki thì mỉm cười hiền như mọi khi.

Mà nói ``bình thường'' thì cũng hơi quá\dots

Ví dụ như chuyện chụp ảnh lớp lúc sáng chẳng hạn.

Tôi vốn là người cầm máy thay cho Uzuki-san, đứng ngoài khung hình là đã ổn, vì tôi vốn là người không thích bị chụp ảnh. Nhưng rồi, cả đám đã nhao nhao kéo tôi vào giữa, ép tôi đứng cạnh hai chị em ấy.

Tôi còn nhớ khoảnh khắc đó chẳng khác nào như ong bị vỡ tổ.

``Kuon-kun, đứng vào đi! Không có cậu thì mất ý nghĩa còn gì!'' cô chị đẩy lưng tôi.

``Onii-sama ở ngoài trông buồn lắm,'' cô em út tiếp lời, nhẹ mà ác.

Tôi còn chưa kịp phản ứng thì Shino-san đã cười tươi, giơ tay ra hiệu: ``Kuon-san đứng ở giữa đi, tớ chừa chỗ rồi này!''

Và thế là tôi bị lôi vào.

Tiếng máy ảnh vang lên tách một cái, giữ nguyên một khoảnh khắc mà đáng lẽ ra tôi không thuộc về, nhưng kỳ lạ là cũng không thấy khó chịu.

Buổi học chiều thì trôi đi nhẹ như gió: làm quen các bạn học mới, chỉ bài cho mấy đứa ngồi bàn cạnh, đổi sổ ghi chép với Sakura-san, nghe Akane-san kể về mấy trò đùa ngớ ngẩn của cô bạn tai mèo của mình.

Tất cả đều thật đến mức tôi không thể không nhận ra: đây là con người thật, không còn là những con rối lặp đi lặp lại vô hồn trong các vòng lặp trước nữa.

Những tiếng cười.

Những câu nói vu vơ.

Cả sự nóng ruột của Akane-san khi bị gọi trả lời.

Mọi thứ đều mang một độ rung cảm mà vòng lặp không bao giờ tái tạo được.

Nhưng\dots\ cũng chính vì vậy mà sự  ``bình thường'' ấy lại khiến tôi bất an. Tôi cảm tưởng rằng mọi thứ vẫn chưa kết thúc.

Và nó đúng như vậy.

Trong giờ nghỉ trưa, tôi ghé qua rừng phía sau trường để thám thính tình hình, và đó cũng là lúc Game-san đã nói với tôi bằng cái giọng trầm trầm đầy cảnh giác hiếm thấy:

``\eng{Thị trấn này chưa hoàn toàn được chữa lành. Ánh sáng trắng chỉ tái tạo phần thân xác, nhưng những vết thương tinh thần, những oán niệm vẫn bám lấy không gian. Cậu đang nhìn thấy dấu vết của họ, không phải ảo ảnh.}''

Tôi không hiểu rõ những gì ông ta nói, chỉ biết rằng\dots\ chúng tôi vẫn chưa xong việc. Không có lỗ xoáy, chưa có đường rời khỏi thế giới này, và rằng chúng tôi sẽ phải ở lại đây một thời gian nữa, ít nhất cho tới khi chúng tôi hoàn thành.

Cái cụm ``hoàn thành'' ấy khiến tôi đăm chiêu suốt cả buổi chiều. Liệu những gì mà Sakura-san nói từ vòng lặp trước đó\dots\ là thật?

Và nếu là thật, thì liệu chúng còn ở đây để tìm cách ám nơi này do những gì mà người dân ở đây đã làm từ thời quá khứ?

\ct

Khi tiếng chuông tan học vang lên, tôi lắc nhẹ đầu để gạt bớt mớ suy nghĩ trong lòng. Thu dọn đồ đạc, tôi cùng Kaigou-chan và Suzuki xách cặp bước ra khỏi lớp.

Thị trấn Aogami lúc chiều tà mang một vẻ đẹp mà tôi chưa từng thấy trong bất cứ vòng lặp nào: nắng nghiêng vàng óng trên các mái ngói, gió nhẹ luồn qua những tán anh đào còn sót cánh, và đâu đó thấp thoáng các bà cô đẩy xe bán hàng, vài đứa trẻ tiểu học chạy nô đùa dưới bóng đèn đường mới bật. Không khí ấy bình yên đến mức khiến người ta quên mất rằng nơi này từng bị nghiền nát bởi lời nguyền không biết bao nhiêu lần.

Ra đến cổng trường, chúng tôi gặp bốn cô gái cùng lớp tôi đang đứng nói chuyện một cách rôm rả.

Shino-san thấy tôi đầu tiên, vẫy tay thật rộng: ``Kuon-san, Kaigou-san, Suzuki-san, các cậu muốn về chung không?''

Kaigou-chan chống hông: ``Đi chứ! Bốn người các cậu đi chung với nhau như thế, không biết nhà mấy cậu hướng nào nhỉ?''

Akane-chan cười toe: ``Nhà bọn mình cũng hướng về phố chính á! Đi chung đi cho vui!''

Sakura-san khoanh tay, nheo mắt nhìn tôi: ``Kuon-san, cậu mà bỏ đi hướng khác là tôi biết liền đấy. Không trốn được đâu.''

``Ơ\dots\ tôi có trốn gì đâu\dots'' tôi bật cười, khoát tay. ``Đi thì đi, đứng đây thêm nắng chói mắt quá.''

Yuka-san thì chỉ khẽ gật đầu, mắt lim dim: ``Hưm\~{} đi thôi\dots'' Rồi cô nhỏ lẩm bẩm cái gì đó nghe như: ``\hl{Đường phía tây hôm nay\dots{} có vẻ ổn hơn nha\~{}.}''

Không ai hiểu cô ấy nói gì. Tôi nhớ cô ấy là một nhà ngoại cảm, vậy nên hẳn Yuka-san nói vậy là có ý gì đó\dots

Cô bạn bờm chó đánh nhẹ vào vai Yuka-san: ``Bà đừng nói mấy câu ghê ghê vậy chứ!''

``Ồ\~{} xin lỗi\~{}'' cô bạn tai mèo đáp, giọng vẫn ngái ngủ.

Thế là chúng tôi đi chung thành một nhóm bảy người, đông đến mức nhìn từ xa như một câu lạc bộ nào đó tụ họp cuối ngày.

Không biết từ lúc nào, đôi chân của chúng tôi tự nhiên chọn đúng con đường cũ, con đường từng xuất hiện trong vô số vòng lặp mà chẳng ai khác ngoài những người đã tái sinh cả thị trấn này như chúng tôi còn nhớ.

Tôi đi ở chính giữa, nhìn những chiếc bóng kéo dài trên mặt đường gạch. Ánh nắng chiều tô viền vàng lên tóc từng người, gió thổi qua mang theo hương mộc quen thuộc của thị trấn Aogami: mùi đất, mùi gió biển xa, mùi của một thị trấn cũ kỹ đang cố thở lại sau nhiều lần chết đi sống lại.

Và tôi\dots

Tôi chỉ vừa thấy được một ngày học đầu tiên bình thường.

Nhưng cái bóng của những điều bất thường vẫn lảng vảng đâu đó quanh tôi, nhất là sau câu nói của Game trong rừng.

Tôi hít một hơi thật sâu.

Hy vọng rằng\dots\ chuyến đi bộ này sẽ diễn ra suôn sẻ.

\ct

Tôi đi đầu cùng Kaigou-chan và Shino-san, những người còn lại bám theo phía sau. Aogami buổi sáng phủ một lớp sương mỏng, vừa đẹp vừa lạnh như thể che giấu điều gì đó. Dù là lần đầu chúng tôi đặt chân đến đây, tôi vẫn cảm nhận rõ có thứ gì đang lệch nhịp trong thị trấn này. Như thể nó vẫn còn bị kéo giật từ vòng lặp cũ về.

Đến bờ sông, nơi mặt nước lặng như tờ, cô chị lớn bỗng dừng lại. Tôi nghe tiếng cô thở dài rất khẽ, giống như thể ký ức nào đó va vào tâm trí cô.

Shino chạm nhẹ vai cô: ``Cậu nhớ gì à?''

Kaigou-chan khẽ gật, nheo mắt nhìn theo dòng nước, nơi cô từng là Saikyou, chưa khôi phục ký ức, từng đứng cùng Touka-san vào một buổi chiều mát rượi.

``Chỉ là cảm cúm thôi.''

Lời nói đó vang lại trong đầu tôi. Cách cô ấy trừng mắt, giọng đổi khác một cách dị thường, theo những gì cô ấy nhớ. Đó là lúc mà cô hỏi tại sao nguyên cả lớp bên cạnh lại vắng mặt, và Touka-san lúc đó đã né tránh. Theo một cách đáng sợ.

Trong hiện tại, bờ sông tối hơn tưởng tượng. Rêu mọc nhiều, đường đi bị đất đá xô lệch. Không có học sinh cười đùa, không có tiếng nước chảy róc rách như trong vòng lặp. Chỉ có sự im lặng.

Akane-san nhìn quanh rồi hỏi: ``Ơ, ở đây trông\dots\ buồn thế? Tưởng thị trấn nhỏ thì phải nhiều người đi bộ hơn chứ?''

Yuka-san thì lại khác. Cô nhìn thẳng xuống mặt nước, mắt hơi trợn: ``\dots Có ai đó đang nhìn tụi mình dưới đó.''

Cả Sakura-san lẫn Akane-san đều quay sang: ``Hả!? Y-Yuka(-san)!? Đừng dọa vậy chứ\dots''

Kaigou-chan rùng mình: ``Lại cái kiểu cảm giác kỳ lạ nữa\dots''

Cô nàng tóc nâu liếc cô bạn tai mèo rất nhanh, đầy cảnh giác, như đang tự hỏi liệu cô ấy có nhìn thấy mảnh sót của vòng lặp trước không.

Tôi không xen vào. Tôi chỉ lặng lẽ nhìn mặt nước. Tôi biết nơi này từng là điểm mốc đầu tiên của vòng lặp, nhưng trong hiện thực, nó tĩnh lặng quá mức, như thể chờ một ai đó khuấy động.

Và nó khiến tôi thấy không đúng.

Đi qua bờ sông thơ mộng nhưng yên tĩnh tới mức kỳ lạ ấy, khu phố truyền thống hiện ra sau vài phút đi bộ. Các cửa hàng đa phần đã đóng hoặc ngừng hoạt động khi trời đã ngả vàng. Gió len vào những tấm biển treo, phát ra âm thanh khô khốc.

Suzuki nhìn quanh rồi cười nhạt: ``Ở vòng lặp, khu này vui hơn nhiều.''

Kaigou-chan gật đầu: ``Ừ. Lúc đó, chúng ta đã gặp mấy cô gái ở đây. Yuka-san, Akane-san, rồi Kaoru-san, Miku-san và Saya-san.''

Shino-san bổ sung thêm: ``Họ đến khu này khi vòng lặp ga tàu bắt đầu. Nếu mình nhớ không nhầm.''

Cô nàng tai chó ngẩng đầu: ``Hể, bọn tôi từng đến đây á? Nghe giống du ngoạn ấy.''

``Không đâu,'' tôi nói, ``lúc đó trông hai người\dots\ đáng sợ thì đúng hơn.''

Yuka-san lại quay về kiểu nhìn không chớp mắt của cô, trầm giọng: ``Ở đây\dots\ không giống như những gì tôi nhớ.''

Akane-san giật mình: ``Hả? Nhưng bọn mình có tới đây bao giờ đâu mà nhớ?''

Cô bạn bờm mèo không trả lời. Cô sờ tay vào bức tường gỗ cũ, rồi hạ giọng: ``Tôi nghe thấy tiếng bước chân. Năm người. Giống bọn mình. Nhưng mà\dots\ không phải bọn mình.''

Sakura-san toát mồ hôi hột, nhìn cô nàng tóc tím đang cảm nhận bằng năng lực của mình: ``Yuka trông đáng sợ thật đấy\dots\ Bộ cậu nhìn thấy ma hả?''

Còn tôi thì suy tính. Nếu Yuka thật sự có thể cảm nhận dư chấn vòng lặp\dots\ thì những gì mà Game-san nói lúc đó có khi lại có cơ sở.

Suzuki đứng cạnh tôi thì thầm:

``Onii-sama, em bắt đầu thấy thị trấn này\dots\ có gì đó kỳ lạ. Dù chỉ là chiều muộn thôi, em vẫn thấy có gì đó\dots''

``Không chỉ em. Anh cũng thấy thế,'' tôi lên tiếng an ủi.

Khi đến gần nhà ga, tim tôi bất giác đập nhanh hơn. Nhà ga Aogami trong vòng lặp từng là tâm điểm, nơi mọi thứ trở nên sai lệch, nơi Yuka-san biến mất giữa chuyến tàu hơi nước lạc thời đại, nơi Akane-san chạy khắp toa la gọi tên bạn mình, và nơi mà hai bản thể của Shino đụng mặt nhau, từ đó dẫn chúng tôi tới gần hơn đến sự thật.

Nhà ga thật bề ngoài trông thì bình thường, thậm chí hơi cũ. Không có hơi nước, không có âm thanh kỳ dị, không có bất cứ cái gì bị lạc từ thời Chiêu Hòa.

Nhưng khi bước lên sân ga, tôi cảm nhận rõ: không khí nơi này vẫn lệch nhịp. Như thể vòng lặp từng đục khoét nơi này và để lại sẹo.

Yuka-san, một lần nữa, là người phản ứng sớm nhất. Cô bước lùi lại, mắt mở to: ``Chỗ này\dots\ đang chạy đi. Có cái gì đó\dots\ đang đi vòng vòng quanh trụ đồng hồ.''

Akane-san vội bám lấy tay cô, toàn thân như đang lạnh sống lưng: ``Yuka! Không có ai hết á!!''

Suzuki run run: ``Yuka-san\dots? Cậu thấy gì vậy?''

Cô nàng ngoại cảm không trả lời. Nhưng gương mặt cô hơi căng ra như thể nhìn thấy đúng cái thứ mà từng kéo bạn mình ra khỏi thực tại trên con tàu hơi nước.

Shino-san tiến đến, giọng trầm và bình tĩnh một cách đáng sợ: ``Yuka-san. Cậu đang nhìn thấy bóng sót của vòng lặp cũ. Đúng không?''

``T-Trụ đồng hồ đang lệch. Nó rung nhẹ. Và có một người đang đi quanh nó. Nhưng người đó\dots\ không có nét mặt.''

Cả nhóm im lặng.

Kaigou khẽ quay sang tôi, thì thầm: ``Cậu cảm nhận được không?''

Tôi hít một hơi thật sâu. ``Có. Nó giống như một lớp ký ức cũ đang trồi lên. Dù vòng lặp biến mất rồi, nhưng\dots\ dấu tích của nó vẫn còn ở đây.''

Sakura-san thêm vào: ``Như thể nhà ga này có thứ gì đó\dots\ không chịu chết.''

Cô nàng tóc nâu gật đầu. Ánh mắt cô sắc lại: ``Đúng. Chúng ta càng đi sâu vào thị trấn, dư ảnh vòng lặp càng dày đặc.''

Còn tôi thì lại cảm thấy rõ một điều: Aogami không bình thường. Dù nó đã được tái sinh lại, nhưng dường như dư chấn của những gì diễn ra từ quá khứ vẫn còn vương vấn tại nơi đây.

Và dường như\dots\ nó vừa thức dậy.

\ct

Chúng tôi rời khỏi ga tàu, bước xuống những bậc cấp đá mòn nhẵn dưới chân. Con đường đất rộng, hai bên là dãy nhà gỗ xám xịt, mái ngói đen bóng nước mưa đọng lại thành vũng nhỏ phản chiếu ánh hoàng hôn. Cổng hàng rào sắt rỉ đỏ lòm nhô ra, níu lấy tà áo kẻ sọc của đứa trẻ chạy qua; tiếng khóa cổng leng keng như thứ âm thanh duy nhất còn sót lại của buổi chiều. Những con hẻm hẹp len lỏi giữa các dãy nhà, chỗ thì chỉ rộng bằng hai vai, tường gạch không vữa, rêu phong bám đầy vết nứt như mạng nhện giăng ngang khu trung tâm. Đâu đó, mùi gỗ tạp chưa ráo hơi sơn cũ lẫn mùi thuốc lá khét của ông lão ngồi trước hiên, khói xám bay lên tan trong gió. Một chú mèo vàng lông xù băng qua đầu chúng tôi, nhảy lên nóc tường, làm rơi xuống vài mảnh vôi trắng vỡ vụn. Ánh nắng xế chiều đỏ au, xuyên qua khe lá khô rơi rụng, in lên mặt đường những hình thoi nhỏ như mắt cáo đang dõi theo từng bước chân.

Tôi định bước tiếp thì Yuka-san đột nhiên khựng lại. Cô đứng yên, đầu hơi nghiêng sang trái như đang nghe thứ gì đó mà người thường không nghe được.

Kaigou-chan hỏi nhỏ: ``Yuka-san? Sao vậy?''

``Không có gì đáng sợ đâu,'' Yuka-san đáp, giọng vẫn đều đều như mọi khi. ``Nhưng mà\dots\ tôi thấy có mấy đứa trẻ.''

Akane-san ngay lập tức giật mình: ``\textbf{Đứa trẻ nào!? Ở đâu!? Tôi không thấy gì hết!}''

CÔ bạn tai mèo chỉ nhẹ nhàng đưa tay, chỉ vào khoảng không giữa hai căn nhà cổ: ``Ở đó. Ba đứa. Chắc tầm tiểu học. Chúng nhìn sang bên này\dots\ rồi lại chạy mất.''

Tôi nheo mắt nhìn theo hướng cô ấy chỉ. Ban đầu chẳng có gì, cho đến khi một làn gió thổi ngược, hất tung lớp bụi trên mặt đất.

Ngay khoảnh khắc đó\dots\ tôi thấy \emph{chúng}.

Những bóng nhỏ. Mờ mờ. Trắng nhạt như sương. Không có mặt, không có tiếng bước chân. Chỉ đứng đó nhìn. Tò mò, lặng lẽ, rồi biến mất sau bức tường gỗ.

Suzuki lập tức bám vào tay áo tôi: ``Onii-san\dots\ cái đó\dots\ cái đó chắc chắn không phải người sống, đúng không!? Em không thích mấy thứ mờ mờ như vậy đâu đấy!?''

Trong khi Yuka-san chỉ nhướng mày một chút, biểu cảm dạng ``à, thì ra là vậy''.

Shino-san trái lại, chỉ nhắm mắt trong vài giây, đầu hơi cúi xuống. Tôi nghe thấy cô khẽ thở ra một hơi nặng trĩu. Cảm giác như cô hiểu quá rõ những thứ đang hiện ra trước mặt.

Còn Sakura-san lại siết chặt hai bàn tay đến mức khớp tay trắng bệch. Nỗi sợ hiện rõ, nhưng xen lẫn với một thứ gì đó khác: một sự thương xót.

``\dots Không lẽ nào\dots\ chúng là những đứa trẻ bị hiến tế sao?'' cô thì thầm, giọng run run. ``Vậy những gì tôi\dots\ tìm hiểu lúc đó là thật?''

Câu nói ấy khiến cả nhóm im bặt. Không khí bỗng như bị bóp nghẹt.

Và rồi ngay lúc ấy, giọng Game vang lên trong đồng hồ tôi. Không còn vẻ đùa cợt, mà trầm, khàn và cẩn trọng hơn nhiều.

``Quả đúng như những gì tôi nghĩ. Có vài linh hồn vẫn còn vướng mắc ở nơi này.''

Tôi siết chặt nắm tay theo bản năng.

Kaigou bước sát lại tôi, thì thầm: ``Cậu cũng thấy chúng\dots\ đúng không?''

``\dots Ừ.'' Tôi đáp. ``Những linh hồn trẻ em đó. Những nạn nhân của hiến tế.''

Khi tôi nhìn kỹ lại, quả thật có hai, ba bóng dáng khác lẩn khuất giữa những ngôi nhà cũ, như đang tìm kiếm điều gì đó\dots\ hoặc ai đó.

Không có ác ý. Không có oán khí.

Chỉ là sự lặng lẽ. Sự cô độc. Và một nỗi buồn không gọi thành lời.

Suzuki run lên: ``Em không chịu được mấy thứ buồn buồn kiểu này, nhìn như đang chờ ai tới đón vậy\dots''

Yuka-san chỉ gật nhẹ: ``Chúng không hại ai đâu. Chỉ là chúng\dots\ bị kẹt lại thôi.''

Akane-san quay trái quay phải như một đứa trẻ chưa hiểu chuyện, dẫu rằng nét mặt cô có chút hoang mang: ``Ủa nhưng\dots\ kẹt lại là sao? Sao lại có linh hồn? Rồi tại sao tôi không thấy gì hết ngoài mọi người vậy?''

Tôi chưa kịp trả lời thì Game-san nói tiếp: ``Không sai. Nhưng điều kỳ lạ là, đáng ra chúng phải được giải thoát khi dòng chảy sinh mệnh của thị trấn được khôi phục. Vậy mà chúng vẫn còn ở đây.''

Chúng tôi nhìn quanh. Gió chiều thổi qua những mái ngói cũ. Chuông gió kêu leng keng trên hiên nhà, nghe như tiếng khóc bị bóp nghẹt.

Tôi cúi xuống, nhìn nền đất sẫm màu dưới chân mình.

``\dots Tức là thị trấn này chưa được khôi phục hoàn toàn.''

Không ai nói gì.

Shino-san chỉ khẽ nhắm mắt, hàng mi run nhẹ, cô hiểu ý nghĩa sâu xa trong câu đó hơn bất kỳ ai.

Kaigou-chan nắm chặt tay áo của tôi.

Yuka-san nhìn theo một bóng hình đang tan vào tường gỗ, ánh mắt nghiêm túc hiếm thấy.

Suzuki rùng mình một cái rồi lùi lại nửa bước.

Sakura-san cắn môi, như sợ mình sẽ khóc.

Còn Akane-san chỉ đứng giữa họ, ngơ ngác và hoang mang: ``Ơ\dots\ mọi người đang nói chuyện\dots\ thật luôn hả?''

Một buổi chiều yên bình ở Aogami, nhưng bên dưới nó là thứ gì đó bất ổn.

Như thể thế giới hiện tại chỉ là lớp da mỏng phủ lên một thực tại cũ đang cố trồi lên khỏi mặt đất.

\ct

Mặt trời gần khuất sau dãy núi. Sắc cam nhạt chuyển dần sang tím, phủ lên thị trấn một màu u ám dị thường.

Chúng tôi đứng giữa lối đi lát đá, lặng im như thể sợ rằng tiếng động nhỏ nhất có thể quấy rầy những linh hồn đang len lỏi quanh đây.

Shino-san mở mắt, nhìn thẳng vào khoảng không phía trước, nơi một bóng nhỏ vừa biến mất.

``Những linh hồn trẻ em đáng ra không nên ở đây\dots'' cô thì thầm. ``Nếu vòng lặp bị phá vỡ hoàn toàn\dots\ đúng ra mọi thứ phải được đưa về đúng trật tự.''

``Nhưng trật tự ở đây\dots\ đã bị bẻ cong,'' tôi nói.

Kaigou nhìn quanh, sắc mặt nặng nề: ``Giống như thị trấn này\dots\ vẫn nhớ vòng lặp. Như thể quá khứ của nó không chịu qua đi.''

Suzuki lùi sát lại tôi hơn: ``Em bắt đầu sợ nơi này rồi đó, Onii-sama\dots\ Chúng ta phải làm gì bây giờ?''

Yuka-san lại ngược hẳn. Cô chỉ nghiêng đầu quan sát như thể đang cố lắng nghe một giai điệu rất xa: ``Không buồn, nhưng rất lạnh. Giống như mấy đứa trẻ muốn nói gì đó\dots\ nhưng không thể nói được.''

Akane-san lập tức hoang mang cực độ, như thể bị gạt ra khỏi cuộc trò chuyện: ``Ủa!? Làm sao mà mọi người biết được nhiều thứ vậy?''

Sakura-san đáp thay mọi người, giọng có chút nghiêm trọng: ``Vì chúng tôi đã từng ở đây. Một phiên bản khác của nơi này.''

Tôi khẽ nhắm mắt khi nghe câu đó.

Phiên bản khác. Vòng lặp. Những cái chết. Những đứa trẻ bị hiến tế.

Tất cả\dots\ vẫn chưa hoàn toàn biến mất.

Game lại lên tiếng, lần này chậm rãi hơn, như đang cân nhắc từng chữ: ``Tàn dư không phải là vấn đề. Điều đáng nói là: tại sao dòng chảy sinh mệnh đã hồi phục mà những linh hồn này vẫn chưa được giải thoát.''

Tôi ngẩng đầu lên.

Câu hỏi đó không ai trong chúng tôi có câu trả lời.

Gió thổi qua, lạnh buốt.

Chuông gió lại rung lên khe khẽ.

Những bóng dáng nhỏ kia lại xuất hiện thoáng qua rồi biến mất, như thể nhắc nhở chúng tôi rằng thị trấn này chưa bao giờ thật sự yên ổn.

Và trong khoảnh khắc ngắn ngủi khi hoàng hôn chạm đáy, tôi cảm thấy rất rõ: Aogami, nơi thực tại duy nhất chúng tôi đang sống, vẫn đang gợn sóng bởi quá khứ chưa khép lại.

Chúng tôi vẫn đứng đó vài giây sau khi những bóng dáng nhỏ cuối cùng tan vào bức tường gỗ mục. Không một tiếng động, không một tiếng thì thầm, như thể cả dãy nhà đang nín thở cùng chúng tôi.

Rồi, Kaigou-chan lên tiếng khẽ: ``\dots Đi thôi. Ở đây lâu\dots\ khó chịu lắm.''

Không một ai lên tiếng phản đối, khẽ gật đầu.

Chúng tôi bước tiếp về phía con đường chính của thị trấn.

Và thật kỳ lạ, chỉ cách đó vài dãy nhà thôi, không khí đã khác hẳn.

Tiếng người nói ríu rít từ những quán ven đường, tiếng xe đạp chạy ngang, tiếng chó sủa xa xa. Hương bánh nướng từ tiệm gần góc phố hòa vào mùi gió chiều. Aogami ở đây nhộn nhịp, sống động, bình thường đến mức khó tin rằng chỉ vài phút trước chúng tôi vừa đứng giữa một nơi lạnh lẽo đến nghẹt thở.

Suzuki bước nhanh hơn một chút, cố lấp đầy khoảng lặng ban nãy: ``May quá\dots\ cuối cùng cũng ra chỗ có người. Em thề là tim em nãy giờ đập to đến mức tưởng ai cũng nghe thấy\dots''

Shino-san không nói gì. Cô nắm hai tay vào nhau, ánh mắt hướng xuống mặt đường lát đá, gương mặt trầm lại.

Akane-san thì vẫn\dots\ là Akane-san. Cô nhìn quanh, rồi thở phào: ``Trời ơi, nãy vẫn chưa hiểu gì hết\dots\ giờ thấy người qua lại bình thường là đỡ lo hẳn rồi đó\~{}''

Tôi nghĩ cái đỡ lo hơn hẳn ấy sẽ giảm xuống nhiều hơn nếu cô thật sự nhìn thấy những gì Yuka-san thấy, nhưng thôi, chuyện đó chắc để sau.

Thị trấn bình yên này giống như đang cố che giấu một cái gì đó dưới lớp vỏ của nó. Giống như cách mà các vòng lặp trước đã thực hiện, nhưng theo một cách khác.

Và tất cả chúng tôi đều cảm nhận được điều ấy, theo những hướng khác nhau.

Shino là người phá tan sự im lặng. Giọng cô nhỏ, nhưng rõ ràng: ``\dots Nếu chúng vẫn còn ở lại\dots\ có lẽ, đâu đó vẫn còn điều chưa được tha thứ.''

Tôi dừng bước một chút.

Những dải mây tím loang trên bầu trời như những vệt mực kéo dài vô tận. Tôi ngẩng lên, hít một hơi dài, cố gỡ cảm giác nặng nề đang đè lên ngực mình.

``Có lẽ,'' tôi nói, ``hành trình lần này sẽ không chỉ là sống lại cuộc đời học sinh\dots\ mà còn là giúp những linh hồn đó được yên nghỉ.''

Một cơn gió chiều nhẹ lướt qua, mang theo mùi gỗ cũ và nắng tàn. Đồng hồ trên cổ tay tôi nhấp nháy một ánh sáng nho4, như phản hồi lại lời tôi vừa nói.

Ngay sau đó, giọng Game vang lên, trầm và nghiêm túc đến mức mọi âm thanh xung quanh như chìm xuống: ``Khôi phục thật sự\dots\ mới chỉ bắt đầu thôi, Kuon-user.''

Ánh sáng ấy tắt ngay lập tức, để lại sự im lặng mơ hồ khó diễn tả.

Cuối phố, chúng tôi dừng lại ở ngã ba, nơi mà mỗi người sẽ tách ra để về nhà trọ của mình.

Sakura-san vẫy tay nhẹ: ``Vậy\dots\ gặp lại mọi người vào ngày mai nhé.''

Akane-san giơ cả hai tay lên như đang cổ vũ: ``Ngày đầu trôi qua ổn rồi! Mai cố lên nữa nha!''

Shino-san mỉm cười. Một nụ cười thật, đẹp và nhẹ hơn hẳn nỗi lo ban nãy: ``Hẹn gặp lại mọi người.''

Kaigou chìa nắm đấm nhỏ về phía tôi: ``Cậu đừng đăm chiêu quá nha, Kuon-kun. Khuya lại mất ngủ ráng chịu đó.''

Suzuki đứng sát tôi, nói nhỏ: ``Onii-sama nhớ khóa cửa phòng nhé\dots\ em bắt đầu\dots\ sợ mấy chỗ tối rồi đó\dots''

Còn Yuka-san chỉ cúi đầu thật nhẹ, chậm rãi nói: ``Ngày mai\dots\ có thể chúng sẽ lại xuất hiện. Mọi người cảnh giác.''

Chúng tôi chia nhau ra, mỗi người bước về một hướng. Âm thanh của thị trấn lại trở về với nhịp sống quen thuộc, như thể không có gì xảy ra.

Tôi đứng đó thêm một lúc, trên tay vẫn còn cảm giác lạnh từ cơn gió khi nãy. Ngẩng lên nhìn bầu trời tím nhạt, trải dài như một biển mực loang.

``Onii-sama,'' Suzuki khẽ kéo tay áo tôi, ``mấy đứa trẻ đó\dots\ chúng có cần chúng ta giúp không?''

Tôi chưa kịp trả lời, Kaigou-chan đã cất tiếng, giọng trầm hơn mọi khi: ``Không phải giúp đỡ thông thường đâu. Chúng bị kẹt giữa hai lớp thời gian.''

``Vậy\dots\ chúng ta phải làm gì?'' cô bé siết chặt tay tôi.

Tôi thở dài: ``Tìm ra cách đưa họ về đúng chỗ. Trước hết, cần biết họ bị kẹt vì điều gì.''

Trong đôi mắt tôi phản chiếu sắc tím ấy, một chút bình thản, một chút mệt mỏi, và một quyết tâm mơ hồ nhưng rất thật dành cho những ngày sắp tới.

Aogami\dots\ vẫn chưa hoàn toàn bình ổn.

Và tôi biết, hành trình khám phá và khôi phục hoàn toàn thị trấn này mới chỉ vừa bắt đầu.

\end{document}